
\medskip
\phantomsection\label{cor:sga3-VIB-5.5}%
\noindent\textbf{Corollary 5.5.---}
\emph{Let \(S\) be a scheme, and let \(G\) be an \(S\)-group
scheme, locally of finite presentation, with connected fibers, and
universally open over \(S\). Then \(G\) is separated and of finite
presentation over \(S\).}%

Indeed, by \SGARefLink{cor:sga3-VIB-3.6}{3.6} and \SGARefLink{cor:sga3-VIB-5.4}{5.4}, \(G\) is quasi-compact and separated over
\(S\). Since \(G\) is locally of finite presentation over \(S\),
it is therefore of finite presentation over \(S\).

\medskip
\phantomsection\label{par:sga3-VIB-5.6}%
\noindent\textbf{5.6.---}
\emph{Proof of \hyperref[thm:sga3-VIB-5.3]{Theorem 5.3}.}
Before proving \SGARefLink{thm:sga3-VIB-5.3}{5.3}, let us establish a few lemmas.

\medskip
\phantomsection\label{lem:sga3-VIB-5.6.0}%
\noindent\textbf{Lemma 5.6.0.---}%
\footnote{\label{note:sga3-VIB-38}Editor's note: we have added this
lemma, communicated by O.\ Gabber. It improves
EGA~IV\(_3\), 14.4.1.2 and corrects the proof of
\emph{loc.\ cit.}, 14.4.1.3 without changing its hypotheses
\textup{(compare the erratum \textup{(ErrIV, 38)} in
EGA~IV\(_4\))}.}
\begin{enumerate}
\item[\textup{\SGAInlineRefTarget{lem:sga3-VIB-5.6.0:item:i}(i)}] \emph{Let \(A\subset B\) be integral domains,
  with \(B\) integral over \(A\), and let
  \(\mathfrak q\in\operatorname{Spec}(B)\) be such that \(A\) is
  unibranch at
  \(\mathfrak p=\mathfrak q\cap A\). Then the morphism}
  \[
    \pi:\operatorname{Spec}(B)\longrightarrow\operatorname{Spec}(A)
  \]
  \emph{is open at \(\mathfrak q\).}
\item[\textup{\SGAInlineRefTarget{lem:sga3-VIB-5.6.0:item:ii}(ii)}] \emph{Let \(X,Y\) be irreducible preschemes,
  let \(f:X\to Y\) be a dominant morphism, and let \(x\in X\) be a
  point at which \(f\) is quasi-finite, such that \(y=f(x)\) is a
  unibranch point of \(Y\). Then \(f\) is open at \(x\). In
  particular, if}
  \[
    Y=\operatorname{Spec}(\mathcal O_{Y,y})
  \]
  \emph{is a local prescheme with closed point \(y\), then
  \(f(U)=Y\) for every neighborhood \(U\) of \(x\).}
\end{enumerate}

\textup{\SGARefLink{lem:sga3-VIB-5.6.0:item:i}{(i)}}
Let \(K\) (respectively, \(L\)) be the fraction field of \(A\)
(respectively, \(B\)), let \(A'\) be the normalization of \(A\),
and let \(B'\) be the subring of \(L\) generated by \(A'\) and
\(B\). Then \(B'\) is integral over \(A'\). Set
\[
  Y=\operatorname{Spec}(A),\quad
  X=\operatorname{Spec}(B),\quad
  Y'=\operatorname{Spec}(A'),\quad
  X'=\operatorname{Spec}(B').
\]
We then have the commutative diagram
\begin{center}
  \phantomsection\label{diag:sga3-VIB-5.6.0-integral-square}%
  \begin{tikzcd}[row sep=large,column sep=large]
    X' \arrow[r,"\pi'"] \arrow[d] &
    Y' \arrow[d,"\phi"'] \\
    X \arrow[r,"\pi"] &
    Y
  \end{tikzcd}
\end{center}
in which all morphisms are integral and surjective.

Since \(A\) is unibranch at \(\mathfrak p\), \(Y'\) has a unique
point \(\mathfrak p'\) above \(\mathfrak p\). Consequently, if
\(U\) is an open neighborhood of \(\mathfrak p'\) in \(Y'\), then
\[
  \phi(Y'-U)
\]
is a closed subset not containing \(\mathfrak p\), so that its
open complement is contained in \(\phi(U)\). This shows that
\(\phi:Y'\to Y\) is open at \(\mathfrak p'\). It is therefore
enough to show that \(\pi'\) is open, and we are reduced to the
case where \(A=A'\) is normal.

Let \(N\) be a quasi-Galois extension of \(K\) containing \(L\),
let
\[
  \widetilde X=\operatorname{Spec}(\widetilde B),
\]
where \(\widetilde B\) is the integral closure of \(B\) in \(N\),
and let \(G=\operatorname{Aut}(N/K)\). Since
\(\widetilde X\to X\) is surjective, it is enough to show that
\[
  \widetilde\pi:\widetilde X\longrightarrow Y
\]
is open. Let \(U\) be an open subset of \(\widetilde X\), and set
\[
  U'=\bigcup_{g\in G}gU.
\]
Because \(G\) acts transitively on the fibers of
\(\widetilde X\to Y\)
(cf. [BAC], Chap.~V, \S2.3, Prop.~6), one has
\[
  \widetilde\pi(U)=\widetilde\pi(U'),
\]
and the latter is the open complement of the closed subset
\[
  \widetilde\pi(\widetilde X-U').
\]
This proves \textup{\SGARefLink{lem:sga3-VIB-5.6.0:item:i}{(i)}}.

\textup{\SGARefLink{lem:sga3-VIB-5.6.0:item:ii}{(ii)}}
We may suppose that \(Y\) and \(X\) are reduced. By Zariski's Main
Theorem (cf. EGA~IV\(_3\), 8.12.9), there exist affine open
neighborhoods \(U\) of \(x\) and \(V=\operatorname{Spec}(A)\) of
\(y\), such that \(f(U)\subset V\), and a factorization
\begin{center}
  \phantomsection\label{diag:sga3-VIB-5.6.0-zariski-factorization}%
  \begin{tikzcd}[row sep=large,column sep=large]
    U
      \arrow[r,hook,"j"]
      \arrow[dr,"f"'] &
    V' \arrow[d,"u"] \\
    {} &
    V
  \end{tikzcd}\,,
\end{center}
where \(j\) is an open immersion and \(u\) is finite. Replacing
\(V'\) by the closure of \(j(U)\), we may suppose that \(V'\) is
irreducible, hence that \(V'=\operatorname{Spec}(B)\), where \(B\)
is an integral \(A\)-algebra, finite over \(A\). Moreover, since
\(f\) is dominant, so is \(u\), and the morphism \(A\to B\) is
therefore injective.

Since \(A\) is unibranch at \(y\), it follows from \textup{\SGARefLink{lem:sga3-VIB-5.6.0:item:i}{(i)}}
that \(u\) is open at \(j(x)\), and hence that
\[
  f=u\circ j
\]
is open at \(x\). This proves the first assertion of
\textup{\SGARefLink{lem:sga3-VIB-5.6.0:item:ii}{(ii)}}. The second follows because, if \(Y\) is a local
prescheme with closed point \(y\), every open subset containing
\(y\) equals \(Y\). In the case
\(Y=\operatorname{Spec}(\mathcal O_{Y,y})\), one may also use,
instead of EGA~IV\(_3\), 8.12.9, the local form of Zariski's Main
Theorem found, for example, in [Pes66] or [Ray70b], Chap.~IV,
Th.~1.
